% related_work.tex — Section 6: Related Work

\section{Related Work}\label{sec:related}

Our work draws on several research threads---multi-agent economics,
financial fraud detection, blockchain analytics, adversarial machine
learning, Web2 bot detection, AI agent security, and the recently active
agentic-commerce-security literature---while occupying a position that
none of them does: the empirical, on-chain \emph{detection} of agent
activity grounded in the structural instability of human behavioral
invariants when autonomous AI agents become the transacting entities.

%% -----------------------------------------------------------------
\subsection{Multi-Agent Economic Systems}\label{subsec:rw-maes}

Agent-based computational economics (ACE) models heterogeneous agents
that interact through market institutions.  Tesfatsion's foundational
ACE framework~\cite{tesfatsion2006} demonstrates how macroscopic
economic phenomena emerge from micro-level agent rules, while
Lowe~\etal\cite{lowe2017} extend multi-agent reinforcement learning
(MARL) to mixed cooperative--competitive settings where agents must
simultaneously collaborate and compete.

On the mechanism-design side, Nisan~\etal\cite{nisan2007} provide the
algorithmic game-theory foundations for strategic allocation and auction
problems, while Babaioff~\etal\cite{babaioff2021} study approximately
optimal mechanisms for additive buyers.  More recently,
Park~\etal\cite{park2023} construct generative agents---LLM-driven
entities that simulate believable human behavior in sandboxed
environments.

A common thread unites all of these lines: agents are modeled as
\emph{participants in} markets designed for (and by) humans.  None
considers the case where agents autonomously \emph{initiate}
transactions, hold wallet keys, and operate at machine speed---precisely
the setting enabled by A2A commerce platforms such as OpenClaw and
Moltbook.  Our threat model (\Cref{sec:threat}) fills this gap by
deriving fraud chains directly from documented platform capabilities.

%% -----------------------------------------------------------------
\subsection{Financial Fraud Detection}\label{subsec:rw-fraud}

Graph-based approaches dominate recent fraud-detection research.
Pourhabibi~\etal\cite{pourhabibi2020} provide a systematic review of
graph-based anomaly detection in payment networks, while
Hilal~\etal\cite{hilal2022} offer a broader survey of anomaly-detection
techniques covering statistical, machine-learning, and deep-learning
methods applied to credit-card fraud, insurance fraud, and money
laundering.  Chandola~\etal\cite{chandola2009} supply the broader
anomaly-detection taxonomy on which many financial applications build,
and Carcillo~\etal\cite{carcillo2021} show how supervised and
unsupervised learning can be combined for credit-card fraud detection.

These methods generally assume that the target distribution is stable
enough for deviations to be meaningful.  In consumer banking, that
distribution is shaped by human behavioral baselines: circadian
transaction patterns, geographic plausibility, spending-category
distributions, and biometric-backed authentication signals.  Our
invariant analysis (\Cref{sec:detection}) studies the narrower failure
mode that occurs when the transacting entity is an AI agent rather than
a human and those baselines are no longer stable.

%% -----------------------------------------------------------------
\subsection{Blockchain Analytics}\label{subsec:rw-blockchain}

Address clustering and de-anonymization have matured considerably since
Meiklejohn~\etal\cite{meiklejohn2016} characterized Bitcoin payment
flows through heuristic clustering.  Chen~\etal\cite{chen2020} extend
graph analysis to Ethereum, exploiting smart-contract call graphs and
token-transfer patterns.  Victor and Weintraud~\cite{victor2021} develop
methods to detect and quantify wash trading on decentralized exchanges,
identifying coordinated self-dealing through temporal and volumetric
anomalies.

Our work extends blockchain analytics in two directions.  First, we
target the specific on-chain footprint of AI agents with documented
platform capabilities (\eg ERC-8004 registrations, ERC-4337 account
abstraction wallets).  Second, we validate our threat-model predictions
against real Base-chain transaction data (\Cref{sec:evaluation}),
bridging the gap between theoretical attack taxonomies and empirical
on-chain evidence.

%% -----------------------------------------------------------------
\subsection{Adversarial ML for Transaction Monitoring}\label{subsec:rw-advml}

A substantial body of work studies adversarial evasion of ML-based
detectors.  Biggio~\etal\cite{biggio2013} formalize evasion attacks
against classifiers at test time, and Biggio and Roli~\cite{biggio2018}
review the broader adversarial-machine-learning trajectory and its
evaluation pitfalls.  In multi-agent reinforcement learning,
Gleave~\etal\cite{gleave2020} show that adversarial policies can attack
victim agents through the environment rather than by directly modifying
inputs.  Earlier, Aleskerov~\etal\cite{aleskerov1997} introduced
\textsc{CardWatch}, a neural-network fraud detector that established
many of the behavioral features still used today.

Our invariant-shift argument differs from this literature in a key
respect.  Adversarial ML attacks perturb \emph{individual inputs} to
cross a decision boundary while remaining close to the original
distribution.  The agent-economy threat is structurally different: agents do
not need to perturb their behavior to evade detection---their
\emph{normal operating behavior} falls outside the distribution on
which detectors were trained.  Conversely, when agents \emph{do}
attempt mimicry (CHAIN\_6), the behavioral mimicry paradox
(\cref{subsec:chain6}) shows that deterministic imitation of a
stochastic process produces a distinct statistical signature that is
\emph{easier} to detect than the original agent behavior.  This
contrasts with the adversarial ML setting, where evasion generally
makes detection harder.

%% -----------------------------------------------------------------
\subsection{Web2 Bot Detection}\label{subsec:rw-botdetect}

The problem of distinguishing automated actors from humans has decades
of history in web security.  Bursztein~\etal\cite{bursztein2014}
demonstrate generic CAPTCHA-solving attacks, showing that
challenge--response mechanisms---the canonical defense against
bots---are increasingly fragile.  Iliou~\etal\cite{iliou2021} combine
web-server logs with mouse behavioral biometrics to detect advanced
web bots, achieving high accuracy through features that overlap
significantly with our Temporal Consistency signal: timing regularity,
interaction cadence, and session-length distributions.
Jonker~\etal\cite{jonker2019} study fingerprint surface-based
detection of bot detectors, documenting the adversarial arms race
between bot operators and detection systems.

The nine invariants identified in our threat model
(\cref{sec:threat}) generalize several assumptions underlying Web2
bot detection: that users exhibit circadian patterns, that
interaction timing is stochastic, and that identity creation carries
nontrivial cost.  All three assumptions fail for AI agents, just as
they have progressively eroded for sophisticated web bots.  The
evasion--counter-evasion dynamics documented in two decades of bot
detection research~\cite{bursztein2014,jonker2019} provide a useful
precedent for anticipating how agent-economy fraud will evolve:
initial reliance on simple behavioral heuristics, followed by an
arms race toward structural and graph-based signals---precisely the
trajectory our detection framework anticipates.

%% -----------------------------------------------------------------
\subsection{AI Agent Security}\label{subsec:rw-aisec}

The AI-security community has focused primarily on compromising agents
from the outside.  Greshake~\etal\cite{greshake2023} demonstrate
indirect prompt injection against LLM-integrated applications, and
Carlini~\etal\cite{carlini2023} show that training-time poisoning can
implant persistent backdoors.  Emerging interoperability standards---the
Agent-to-Agent (A2A) protocol~\cite{a2aprotocol2025} and the Model
Context Protocol (MCP)~\cite{mcp2024}---define how agents discover and
invoke one another's capabilities, yet neither specifies fraud-detection
hooks or transaction-rate governance.

Orthogonally, Sybil attacks have been studied since
Douceur~\cite{douceur2002} formalized the problem, with
Jiang~\etal\cite{jiang2020} proposing reputation-based defenses for
peer-to-peer networks.  Our \textsc{Chain\_4} threat
(Disposable Agent Army, \Cref{sec:threat}) is a Sybil attack executed at
machine speed with near-zero marginal cost: each disposable agent is
created, used for a single fraudulent transaction, and discarded before
reputation systems can respond.

Beyond external compromise of individual agents, the broader systemic
risks of agentic AI---including agent collusion and coordinated
multi-agent behavior---have been surveyed by
Lazer~\etal\cite{lazer2026agentic}, which frames at the security level
the same collective-behavior problem our \textsc{Chain\_7} swarm gap
(\Cref{sec:limitations}) targets at the detection level.  Documented
incidents now exist: in May 2026, a hidden Morse-code prompt injection
delivered through a social-media reply caused an agent to drain roughly
\$150K--\$200K from an agent-controlled wallet on the same Base chain we
study~\cite{grokbankr2026}.  Such incidents are agent-\emph{as-victim}
hijackings---squarely within the indirect-prompt-injection thread of
Greshake~\etal\cite{greshake2023}---rather than the agent-\emph{as-fraudster}
A2A commerce our threat model addresses, but they confirm that
agent-mediated on-chain payment fraud has moved from hypothesis to public
record.

%% -----------------------------------------------------------------
\subsection{Agentic Commerce Security}\label{subsec:rw-agentcommerce}

Concurrent with this work, a systematization literature on
agentic-commerce security has emerged.  Zhang~\etal\cite{zhang2026sok}
systematize blockchain agent-to-agent payments across a
discovery--authorization--execution--accounting lifecycle and identify
failure classes including weak intent binding and misuse under valid
authorization.  Mao~\etal\cite{mao2026sok} build a unified security
framework spanning ERC-8004, AP2, x402, and related standards,
enumerating twelve cross-layer attack vectors and a layered defense
architecture; several of their categories overlap our taxonomy (their
\emph{market manipulation} with our \textsc{Chain\_8}, their
\emph{inter-agent trust} with \textsc{Chain\_4}/\textsc{Chain\_5}).  These
works taxonomize the threat and protocol surface, but neither measures
whether agent activity is distinguishable on real transaction data, nor
defines the behavioral invariants whose failure explains \emph{why}
human-assumption detectors miss agents.  Our contribution is
complementary and detection-side: we operationalize a subset of these
threat classes as five on-chain signals and report three-stage empirical
validation, including an intentionally weak real-data ranking result
(AUC\,=\,0.515) that the taxonomy literature has no mechanism to surface.

On the prevention side, Li~\etal\cite{li2026a402} propose A402, which
binds payment finalization to verified service execution using
TEE-assisted atomic service channels, closing the payment--service
decoupling that plain x402 leaves open.  A402 prevents unauthorized
release by construction at the protocol layer; our framework instead
provides post-hoc behavioral detection over agents and protocols that
lack such guarantees, including the plain-USDC Base population we study.
The two approaches are non-competing and most effective in combination.

%% -----------------------------------------------------------------
\paragraph{Key differentiation.}
To our knowledge, this is the only work that pairs a behavioral-invariant
account of detection failure with empirical on-chain detection metrics on
real agent transactions.  Recent systematizations~\cite{zhang2026sok,mao2026sok}
catalog the threat and protocol surface, and protocol
designs~\cite{li2026a402} harden individual payments, but none reports
detection performance on observed agent activity.  The transition we
anticipated is already underway in 2026: industry guidance now calls
explicitly for ``Know Your Agent'' behavioral
monitoring~\cite{financialbrand2026}, operational telemetry documents
agent-impersonation fraud at scale~\cite{datadome2026}, and risk analyses
report that autonomous agents compress financial-crime
timelines~\cite{trmlabs2026}.  We therefore position this paper not as a
first warning but as an early empirical, detection-side complement to a
now-active agentic-commerce-security literature.
